\chapter*{Zusammenfassung}

Die vorliegende Arbeit befasst sich mit einem Folgeprojekt der vorangestellten Studienarbeit ''Anwendung des Maschinellen Lernens in der bildgebenden Qualitätskontrolle'' und untersucht eine Verbesserung und Visualisierung der Ergebnisse der vorangestellten Studienarbeit.

Zunächst wird nach der zugrunde liegenden Theorie der Verwendung von \ac{API}s das Konzept zur Realisierung der Aufgabenstellung vorgestellt.

Die Umsetzung erfolgt zunächst in einer vom Labor unabhängigen Programmierumgebung und implementiert alle geforderten Aufgabenteile erfolgreich. Nach ersten Tests wird das Programm auf die Laborsysteme übertragen und dort getestet.

Abschließend wird das \ac{CNN} anhand eines neuen Datensatzes evaluiert und mit anderen Architekturen anhand der gängigen Metriken verglichen.

Die Ergebnisse dieser Arbeit liefern Erkenntnisse für die einfache Implementierung einer bidirektionalen Webserver-Applikation und zeigen eine einfache Visualisierungsmöglichkeit der Klassifikationsergebnisse von \ac{CNN}s. Die Ergebnisse bieten Grundlagen für zukünftige Arbeiten im Bereich der Bildverarbeitung und der Qualitätssicherung im Industrieumfeld.


\chapter*{Abstract}

This thesis deals with a follow-up project to the preceding study ''Anwendung des Maschinellen Lernens in der bildgebenden Qualitätskontrolle'' and examines an improvement and visualization of the results of the preceding study.

First, after discussing the underlying theory of using \ac{API}s, the concept for implementing the task is presented.

The implementation is initially carried out in a programming environment independent of the laboratory and successfully implements all required tasks. After initial tests, the program is transferred to the laboratory systems and tested there.

Finally, the \ac{CNN} is evaluated using a new dataset and compared with other architectures based on common metrics.

The results of this work provide insights into the simple implementation of a bidirectional web server application and demonstrate an easy way to visualize the classification results of \ac{CNN}s. The findings lay the groundwork for future work in the field of image processing and quality assurance in an industrial environment.